%%%%%%%%%%%%%%%%%%%%%%%%%%%%%%%%%%%%%%%%%%%%%%%%%%%%%%%%%%%%%%%%%%%%%%%%%%%%%%
%  Unit IV, second half: Mobile Marketing (Chapters 6-8)
%%%%%%%%%%%%%%%%%%%%%%%%%%%%%%%%%%%%%%%%%%%%%%%%%%%%%%%%%%%%%%%%%%%%%%%%%%%%%%

\chapter{Mobile Advertising and Mobile Marketing}

\syllabusstrip{Mobile Marketing: Difference between mobile advertising and
marketing.}

\begin{learningoutcomes}
\begin{itemize}
\item What mobile marketing is, and the scale of mobile media consumption.
\item The difference between mobile advertising and mobile marketing --- the
      single most repeated question in this unit.
\item Mobile web marketing against mobile app marketing.
\item The benefits of mobile marketing, and its five real limitations.
\item The seven characteristics of a successful mobile campaign.
\item Geo-targeting and geo-fencing.
\end{itemize}
\end{learningoutcomes}

\section{What Mobile Marketing Is}

\begin{definitionbox}[Mobile Marketing]
\keytermas{Mobile marketing}{mobile marketing} is a multi-channel digital
marketing strategy aimed at reaching a target audience on their smartphones,
tablets and other mobile devices --- through mobile websites, e-mail, SMS and
MMS, social media, apps and push notifications.
\end{definitionbox}

Mobile is disrupting the way people engage with brands. Everything that can be
done on a desktop is now available on a mobile device --- from opening an e-mail
to visiting a website to reading content, all through a small screen. About
\textbf{80\%} of internet users own a smartphone, and mobile platforms host up to
\textbf{60\%} of digital media consumption.

\section{Mobile Advertising Against Mobile Marketing}

\begin{keybox}[The Most Repeated Question in Unit IV]
The same distinction, set in three different phrasings.
\pyq{CIE-III, Jun 2026; Improvement CIE, 4 Jun 2025; Improvement CIE, 28 Aug 2024 --- 2 M each}
\end{keybox}

\begin{exambox}[The One-Line Answer]
Mobile \textbf{advertising} focuses mainly on paid promotional messages delivered
to mobile devices. Mobile \textbf{marketing} includes overall customer engagement
through SMS, apps, social media and mobile websites.

\smallskip
In one line: \textbf{mobile advertising is a paid subset of the broader mobile
marketing strategy.}
\end{exambox}

\begin{figure}[htbp]
\centering
\begin{tikzpicture}[font=\sffamily\footnotesize]
  % outer set: mobile marketing
  \draw[ocre,line width=1.1pt,rounded corners=7pt,fill=ocrepale]
    (0,0) rectangle (11.2,4.35);
  \node[anchor=north west,text=ocredark,font=\sffamily\bfseries\small]
    at (0.3,4.18) {MOBILE MARKETING --- the whole strategy};
  \node[anchor=north west,text=slate,font=\sffamily\scriptsize,align=left,
        text width=4.4cm] at (0.3,3.55)
    {apps \ \textbullet\ \ mobile websites\\ SMS and MMS \ \textbullet\ \ push
     notifications\\ mobile social \ \textbullet\ \ mobile e-mail\\ QR codes \
     \textbullet\ \ location-based\\ mobile commerce};
  \node[anchor=north west,text=slatelight,font=\sffamily\scriptsize\itshape,
        align=left,text width=4.4cm] at (0.3,1.5)
    {paid \emph{and} organic\\ continuous, two-way\\ measured across the whole
     funnel};

  % inner subset: mobile advertising
  \draw[deepblue,line width=1.1pt,rounded corners=6pt,fill=deepbluepale]
    (5.55,0.55) rectangle (10.7,3.8);
  \node[anchor=north,text=deepblue,font=\sffamily\bfseries\scriptsize]
    at (8.12,3.65) {MOBILE ADVERTISING};
  \node[anchor=north,text=slatelight,font=\sffamily\scriptsize\itshape]
    at (8.12,3.28) {a paid subset --- one tactic};
  \node[anchor=north,text=slate,font=\sffamily\scriptsize,align=center,
        text width=4.6cm] at (8.12,2.88)
    {banner and display ads \ \textbullet\ \ in-app ads\\ mobile search ads \
     \textbullet\ \ interstitials\\ video ads \ \textbullet\ \ rewarded video};
  \node[anchor=north,text=slatelight,font=\sffamily\scriptsize\itshape,
        align=center,text width=4.6cm] at (8.12,1.55)
    {always paid \ \textbullet\ \ campaign-bound\\ largely one-way\\ measured on
     impressions, clicks, CTR};
\end{tikzpicture}
\caption[Mobile advertising as a subset of mobile marketing]%
        {The containment relation is the answer. Mobile advertising sits
         \emph{inside} mobile marketing: it is one paid tactic within a strategy
         that also runs organic, continuous and two-way activity.}
\label{fig:mobsubset}
\end{figure}

\begin{mmlongtable}{Mobile advertising against mobile marketing}{tab:advsmark}%
{@{}P{0.15\textwidth}P{0.38\textwidth}P{0.40\textwidth}@{}}
\rowcolor{ocre}\thd{Basis} & \thd{Mobile Advertising} & \thd{Mobile Marketing}\\
\midrule
\endfirsthead
\rowcolor{ocre}\thd{Basis} & \thd{Mobile Advertising} & \thd{Mobile Marketing}\\
\midrule
\endhead
\rlab{Scope} & Narrow --- one tactic & Broad --- the whole strategy\\
\rowcolor{tablealt}
\rlab{Nature} & Always paid & Paid and organic\\
\rlab{Channels} & Banner advertisements, in-app advertisements, interstitials,
mobile search advertisements, video advertisements & Apps, mobile websites, SMS
and MMS, push notifications, social, e-mail, QR codes, location-based marketing,
mobile commerce\\
\rowcolor{tablealt}
\rlab{Objective} & Immediate visibility, clicks and conversions & Long-term
engagement, relationship and retention\\
\rlab{Duration} & Campaign-bound & Continuous\\
\rowcolor{tablealt}
\rlab{Interaction} & Largely one-way & Two-way and interactive\\
\rlab{Measurement} & Impressions, clicks, click-through rate, cost per install &
Full-funnel: engagement, retention, lifetime value\\
\end{mmlongtable}

\section{Mobile Web Marketing Against Mobile App Marketing}

\begin{keybox}[High Yield]
``List any two differences between Mobile Web marketing and Mobile App
marketing.''
\pyq{SEE, Oct/Nov 2023 --- 2 M}
The two strongest pairs are \emph{access and reach}, and \emph{push and device
access}.
\end{keybox}

\begin{mmlongtable}{Mobile web against mobile app marketing}{tab:webvsapp}%
{@{}P{0.18\textwidth}P{0.37\textwidth}P{0.38\textwidth}@{}}
\rowcolor{ocre}\thd{Basis} & \thd{Mobile Web Marketing} & \thd{Mobile App Marketing}\\
\midrule
\endfirsthead
\rowcolor{ocre}\thd{Basis} & \thd{Mobile Web Marketing} & \thd{Mobile App Marketing}\\
\midrule
\endhead
\rlab{Access} & Through a browser; no installation required & Requires download
and installation\\
\rowcolor{tablealt}
\rlab{Reach} & Broader --- anyone with a browser and a link & Narrower --- only
those who install\\
\rlab{Discoverability} & Found through search engines (SEO) & Found through app
store optimisation (ASO)\\
\rowcolor{tablealt}
\rlab{Engagement depth} & Lower; session-based & Much higher; repeat usage and
habit formation\\
\rlab{Push notifications} & Limited, via web push & Full push notification
capability\\
\rowcolor{tablealt}
\rlab{Offline use} & Generally unavailable & Possible\\
\rlab{Device features} & Limited access to camera, GPS, contacts & Full access to
device hardware and sensors\\
\rowcolor{tablealt}
\rlab{Cost} & Lower to build and maintain & Higher --- separate platform builds\\
\end{mmlongtable}

\section{Benefits of Mobile Marketing}

\begin{itemize}
\item \sfb{Reach} --- mobile phones reach people about 15\% more than internet
      marketing does, and mobile advertising draws more clicks.
\item \sfb{Accessibility} --- the mobile device is always at hand; mobile
      advertising follows people everywhere.
\item \sfb{Time factor} --- people are available on mobile around the clock, which
      is not possible with desktops.
\item \sfb{Cost} --- mobile advertising costs much less than other types of
      advertising, allowing more advertising for the same spend.
\item \sfb{Personalised} --- the message can be tailored; people find mobile
      messages more intimate than other internet marketing methods.
\item \sfb{Reaches a broader market} --- smartphones and tablets are cheaper,
      smaller and more portable than PCs, so people previously excluded by
      financial, geographical or technological barriers are now online.
\item \sfb{Instantaneous results} --- we always carry our phones, and they are
      usually switched on, so the message is received the moment it is sent.
\item \sfb{Easy to work with} --- creating elements for mobile is simpler and less
      costly than for desktop; promotions such as downloadable coupons can be
      stored and redeemed without an internet connection.
\item \sfb{Convenient to use} --- the small screen limits content, which forces
      creators to keep it basic and simple; simpler content also adapts better
      across mobile platforms.
\item \sfb{Tracking response} --- user response can be tracked almost
      instantaneously, and phone numbers work well as unique identifiers because
      people keep them longer than e-mail addresses, which helps build accurate
      buyer personas.
\item \sfb{Huge viral potential} --- mobile content is easily shared, producing a
      domino effect and extra exposure at no additional cost.
\item \sfb{Mass communication made easy} --- more people own phones than
      computers, and mobile allows geo-targeting through GPS and Bluetooth as well
      as demographic targeting.
\item \sfb{Micro-blogging benefits} --- platforms such as Instagram and X place
      influence in the hands of ordinary people, so anyone can act as an
      influencer.
\item \sfb{Mobile payment} --- secure mobile payment environments let users
      purchase and pay bills without physical currency.
\end{itemize}

\section{Limitations}

\begin{warnbox}[The Five Real Constraints]
\begin{itemize}
\item \sfb{Platforms too diverse} --- mobile devices have no single standard;
      screen sizes vary and platforms use different operating systems and
      browsers, so one campaign for all of them is difficult.
\item \sfb{Privacy issues} --- users want their privacy respected, so marketers
      must offer clear opt-out instructions.
\item \sfb{Navigation on a mobile phone} --- small screens and no mouse make
      navigation difficult even with touch, so many advertisements go untouched
      because reviewing each one is tedious.
\item \sfb{Ad fatigue and intrusiveness} --- push notifications and interstitials
      are easily overused, driving uninstalls.
\item \sfb{Regulatory constraints} --- SMS and push marketing are governed by
      consent regulations, including do-not-disturb registries in India and GDPR
      in the European Union.
\end{itemize}
\end{warnbox}

\section{The Seven Characteristics of a Successful Mobile Campaign}

Derived from the Mobile Marketing Association's interviews with members and its
awards analysis.

\begin{enumerate}
\item \sfb{Think ``mobile first''.} The most successful campaigns are developed
      with a mobile-first mentality; starting from a mobile perspective gives the
      campaign a stronger foundation. Consider mobile the \emph{connective tissue}
      to all media, since it supports and strengthens every channel.
\item \sfb{Leverage multi-screen usage patterns.} Smartphone and tablet users work
      across screens --- watching a match on television while using the app for
      other scores, or posting while watching a show. Brands leverage cross-screen
      usage to capture attention.
\item \sfb{Utilise the full spectrum of mobile tools and applications.} Mobile
      isn't just a mobile website \emph{or} a display campaign \emph{or} an app
      --- it is a mobile website \emph{and} a display campaign \emph{and} an app,
      and much more.
\item \sfb{Integrate the mobile campaign into the traditional campaign.} One
      manufacturer generated 39 million total views by integrating its game-time
      app with a concurrent televised championship campaign. Mobile campaigns
      should never be produced in a silo.
\item \sfb{Create a campaign that works across multiple screens.} Tablet users are
      often at the \emph{top} of the sales funnel\ix{sales funnel} doing in-depth initial research,
      while smartphone users are often at the \emph{bottom}, ready to purchase.
\item \sfb{Leverage every phase of the sales funnel.} Mobile visitors use it for
      search and discovery as well as for purchase and brand connection.
\item \sfb{Test your way to success.} Mobile is digital by nature, making it
      perfect for tracking and measuring. The second half of the battle is sifting
      through those insights and adjusting to improve future performance.
\end{enumerate}

\begin{keybox}[The Closing Idea]
Mobile isn't simply a new sales channel or marketing tool. It is a revolutionary
new medium that is transforming not only the way consumers connect with brands,
but how they make purchases and stay engaged with them.
\end{keybox}

\section{Geo-Targeting}

\begin{definitionbox}[Geo-Targeting]
\keytermas{Geo-targeting}{geo-targeting} in mobile marketing is the delivery of
content, offers or advertisements to users based on their geographic location,
determined through GPS, Wi-Fi, Bluetooth beacons, cell-tower triangulation or IP
address.
\pyq{SEE, Jun/Jul 2025 --- 2 M}
\end{definitionbox}

\begin{exambox}[The Role of Geo-Targeting]
It makes messages locally relevant and timely --- sending an offer when the user
is physically near the store (\keytermas{geo-fencing}{geo-fencing}); enabling
``near me'' discovery and click-to-map directions; supporting location-specific
pricing, language and inventory; driving footfall to physical outlets; and
improving ROI by eliminating spend on users outside the serviceable area.
\end{exambox}

\begin{summarybox}
Mobile marketing is a multi-channel strategy reaching users on mobile devices
through websites, e-mail, SMS, social, apps and push notifications; about 80\% of
internet users own a smartphone and mobile carries up to 60\% of digital media
consumption. Mobile advertising is the paid subset of that strategy --- narrow,
always paid, campaign-bound, one-way and measured on clicks --- while mobile
marketing is broad, paid and organic, continuous, two-way and measured across the
whole funnel. Mobile web reaches more people with less depth; apps reach fewer
people with far more depth, full push capability, device access and offline use.
The benefits run from reach, accessibility and round-the-clock availability
through cost, personalisation and viral potential to mobile payment; the
limitations are platform fragmentation, privacy, small-screen navigation, ad
fatigue and consent regulation. Seven characteristics mark a successful campaign,
beginning with mobile-first thinking and ending with continuous testing.
Geo-targeting delivers location-relevant offers and geo-fencing triggers them at a
boundary.
\end{summarybox}

%%%%%%%%%%%%%%%%%%%%%%%%%%%%%%%%%%%%%%%%%%%%%%%%%%%%%%%%%%%%%%%%%%%%%%%%%%%%%%
\chapter{Utilizing Mobile Marketing for Sales Promotions}
%%%%%%%%%%%%%%%%%%%%%%%%%%%%%%%%%%%%%%%%%%%%%%%%%%%%%%%%%%%%%%%%%%%%%%%%%%%%%%

\syllabusstrip{Utilizing mobile marketing for sales promotions.}

\begin{learningoutcomes}
\begin{itemize}
\item Push notifications --- what they are, how they serve sales promotion, and
      the eight strategies that make them work.
\item Ten mobile sales-promotion techniques.
\item The mobile advertising formats, including click-to-call and click-to-map.
\end{itemize}
\end{learningoutcomes}

\section{Push Notifications}

\begin{definitionbox}[Push Notification]
A \keytermas{push notification}{push notification} is a short message sent by an
app directly to a user's device screen, appearing even when the app is not open.
\end{definitionbox}

\begin{keybox}[High Yield]
``Discuss briefly on Push Notification in Mobile Marketing. Explain how
businesses can use them to engage and retain Mobile App users? Provide examples
of effective push notification strategies.''
\pyq{SEE, Oct/Nov 2023, Q8a --- 8 M; Improvement CIE, 28 Aug 2024 --- 2 M}
\end{keybox}

\subsection{How push notifications enhance sales promotions}

\begin{itemize}
\item \sfb{Immediacy} --- a flash sale can be announced to the entire installed
      base within seconds.
\item \sfb{Urgency} --- countdown and limited-stock messaging drives immediate
      action.
\item \sfb{Personalisation} --- notifications triggered by browsing or purchase
      history: ``the jacket you viewed is now 30\% off''.
\item \sfb{Cart recovery} --- abandoned-cart reminders sent minutes after
      abandonment.
\item \sfb{Geo-triggered offers} --- a coupon delivered when the user enters a
      defined radius around a store.
\item \sfb{Re-engagement} --- win-back messages to users who have not opened the
      app in 30 days.
\item \sfb{Loyalty and replenishment} --- points-expiry alerts and reorder
      reminders for consumables.
\end{itemize}

\subsection{Eight effective push notification strategies}

\begin{enumerate}
\item \sfb{Segment before sending} --- never broadcast to the entire base;
      irrelevant pushes drive uninstalls.
\item \sfb{Time it correctly} --- send within the user's active hours in their own
      time zone.
\item \sfb{Cap frequency} --- set an explicit limit per user per week.
\item \sfb{Write with a single clear benefit and a verb}, within roughly 40
      characters for the title.
\item \sfb{Use rich media} --- images and action buttons materially raise
      engagement.
\item \sfb{Deep-link the notification} straight to the relevant screen, never to
      the home screen.
\item \sfb{Ask permission at the right moment} --- request notification consent
      after the user has experienced value, not on first launch.
\item \sfb{A/B test} copy, timing and creative, and measure opt-out rate as
      carefully as click rate.
\end{enumerate}

\begin{warnbox}[The Metric Most Often Ignored]
Push campaigns are usually judged on open rate alone. The \textbf{opt-out rate} is
the more important number, because an opt-out is permanent: a campaign that lifts
opens by 5\% while pushing 2\% of the base to disable notifications has destroyed
future reach to buy present clicks.
\end{warnbox}

\section{Mobile Sales-Promotion Techniques}

\begin{mmlongtable}{Ten mobile sales-promotion techniques}{tab:mobpromo}%
{@{}P{0.24\textwidth}P{0.69\textwidth}@{}}
\rowcolor{ocre}\thd{Technique} & \thd{How it works}\\
\midrule
\endfirsthead
\rowcolor{ocre}\thd{Technique} & \thd{How it works}\\
\midrule
\endhead
\rlab{SMS campaigns} & High open rate, often above 90\%; best for time-critical
offers and verified coupon codes\\
\rowcolor{tablealt}
\rlab{Mobile coupons} & Downloadable or in-app codes redeemable in store or
online; can be used offline\\
\rlab{In-app exclusive offers} & Discounts available only through the app, driving
installs and app engagement\\
\rowcolor{tablealt}
\rlab{Flash sales via push} & Time-boxed offers announced by push to create
urgency\\
\rlab{Location-based offers} & Geo-fenced coupons triggered near a store\\
\rowcolor{tablealt}
\rlab{QR codes} & Bridge print, packaging and outdoor media to a mobile landing
page or offer\\
\rlab{Gamification} & Spin-the-wheel, scratch cards and loyalty streaks to drive
repeat opens\\
\rowcolor{tablealt}
\rlab{Mobile wallet passes} & Offers saved to a device wallet with
location-triggered reminders\\
\rlab{Referral programmes} & One-tap share of a referral code from within the
app\\
\rowcolor{tablealt}
\rlab{Mobile-exclusive early access} & App users get the sale hours before the
website\\
\end{mmlongtable}

\section{Mobile Advertising Formats}

\begin{figure}[htbp]
\centering
\begin{tikzpicture}[font=\sffamily\footnotesize,
    ph/.style={draw=slate!60,line width=1pt,rounded corners=4pt,
               minimum width=2.35cm,minimum height=3.6cm,fill=white},
    fmtlabel/.style={anchor=north,align=center,text width=2.5cm,
                font=\sffamily\bfseries\scriptsize,text=ocredark},
    sub/.style={anchor=north,align=center,text width=2.6cm,
                font=\sffamily\scriptsize,text=slate}]

  % --- click to call
  \node[ph] (p1) at (0,0) {};
  \fill[ocrepale] (-1.05,0.95) rectangle (1.05,1.65);
  \draw[ocre,line width=0.8pt] (-1.05,0.95) rectangle (1.05,1.65);
  \node[font=\sffamily\scriptsize,text=ocredark] at (0,1.3) {Call now};
  \fill[rulegrey] (-1.05,-0.2) rectangle (1.05,0.55);
  \fill[rulegrey] (-1.05,-1.4) rectangle (1.05,-0.6);
  \node[fmtlabel] at (0,-2.05) {CLICK-TO-CALL};
  \node[sub] at (0,-2.5) {a tappable number; lifts CTR by 6--8\%};

  % --- interstitial
  \node[ph] (p2) at (2.95,0) {};
  \fill[deepbluepale] (1.9,-1.55) rectangle (4.0,1.75);
  \draw[deepblue,line width=0.8pt] (1.9,-1.55) rectangle (4.0,1.75);
  \node[font=\sffamily\scriptsize,text=deepblue,align=center] at (2.95,0.4)
    {full-screen\\ ad};
  \node[font=\sffamily\scriptsize,text=slatelight] at (3.78,1.6) {\texttimes};
  \node[fmtlabel] at (2.95,-2.05) {INTERSTITIAL};
  \node[sub] at (2.95,-2.5) {appears within an app; click through or close};

  % --- click to map
  \node[ph] (p3) at (5.90,0) {};
  \fill[greenpale] (4.85,-0.35) rectangle (6.95,1.75);
  \draw[greenok,line width=0.8pt] (4.85,-0.35) rectangle (6.95,1.75);
  \draw[greenok,line width=0.7pt] (5.0,0.5) -- (6.2,1.2) -- (6.8,0.3);
  \fill[warnred] (6.2,1.2) circle (0.09);
  \node[font=\sffamily\scriptsize,text=slate] at (5.90,-0.75) {nearest store};
  \node[fmtlabel] at (5.90,-2.05) {CLICK-TO-MAP};
  \node[sub] at (5.90,-2.5) {geo-targeted; opens a map, then contact details};

  % --- expandable
  \node[ph] (p4) at (8.85,0) {};
  \fill[ocre!25] (7.8,-1.55) rectangle (9.9,1.75);
  \draw[ocre,line width=0.8pt] (7.8,-1.55) rectangle (9.9,1.75);
  \draw[-{Latex[length=1.7mm]},ocredark,line width=0.7pt] (8.3,0.6) -- (7.95,0.95);
  \draw[-{Latex[length=1.7mm]},ocredark,line width=0.7pt] (9.4,0.6) -- (9.75,0.95);
  \draw[-{Latex[length=1.7mm]},ocredark,line width=0.7pt] (8.3,-0.4) -- (7.95,-0.75);
  \draw[-{Latex[length=1.7mm]},ocredark,line width=0.7pt] (9.4,-0.4) -- (9.75,-0.75);
  \node[font=\sffamily\scriptsize,text=slate,align=center] at (8.85,0.1)
    {rich media\\ expands};
  \node[fmtlabel] at (8.85,-2.05) {EXPANDABLE};
  \node[sub] at (8.85,-2.5) {canvas ad expands to cover the whole screen};
\end{tikzpicture}
\caption[Four mobile advertising formats]%
        {Four of the principal formats. Each trades intrusiveness against intent:
         click-to-call and click-to-map serve a user already looking for a
         business, while interstitials and expandables interrupt one who is not.}
\label{fig:mobformats}
\end{figure}

\begin{mmlongtable}{Mobile advertising formats}{tab:mobformats}%
{@{}P{0.24\textwidth}P{0.69\textwidth}@{}}
\rowcolor{ocre}\thd{Format} & \thd{Description}\\
\midrule
\endfirsthead
\rowcolor{ocre}\thd{Format} & \thd{Description}\\
\midrule
\endhead
\rlab{Click-to-Call} & Displays a phone number the user can tap to be connected
directly to a call centre or sales centre. Research shows click-to-call
advertisements drive a 6--8\% average increase in click-through rates\\
\rowcolor{tablealt}
\rlab{Interstitials} & Interactive advertisements that appear within an app; once
the user opens the app the interstitial displays, and the user can click through
to a landing page or close it\\
\rlab{Click-to-Map} & Geo-targeting sends messages to users in a specific
location; clicking the advertisement opens a map identifying the nearest store,
and clicking the map displays the contact information on the smartphone\\
\rowcolor{tablealt}
\rlab{Canvas and expandable advertisements} & When clicked, the advertisement
expands to cover the entire phone screen; it can be animated or incorporate rich
media to enhance the user experience\\
\rlab{Mobile search advertisements} & Paid results on mobile search pages, often
with location extensions\\
\rowcolor{tablealt}
\rlab{In-app banner and native advertisements} & Display units embedded in the
app's content flow\\
\rlab{Rewarded video} & The user watches a video in exchange for in-app currency
or content --- high completion rates\\
\end{mmlongtable}

\begin{keybox}[Do Not Send Mobile Advertising to a Static Page]
Mobile display advertisements are effective for building brand awareness\ix{brand awareness} and
generating clicks, leads and conversions. They need not link to a static landing
page --- linking to a \emph{dynamic} page or experience produces improved
engagement and higher conversions.
\end{keybox}

\begin{summarybox}
A push notification is a short message an app sends to the device screen even when
closed, and it serves sales promotion through immediacy, urgency,
personalisation, cart recovery, geo-triggered offers, re-engagement and loyalty
prompts. Eight strategies make it work: segment, time correctly, cap frequency,
write one benefit with a verb, use rich media, deep-link, ask permission after
value, and A/B test while watching opt-out as closely as opens. Ten promotion
techniques span SMS, coupons, in-app exclusives, push flash sales, geo-fenced
offers, QR codes, gamification, wallet passes, referrals and mobile-exclusive
early access. The advertising formats run from click-to-call, which lifts
click-through by 6--8\%, through interstitials, click-to-map, canvas and
expandable units, mobile search advertisements, in-app banners and rewarded video
--- and all of them convert better into a dynamic destination than a static one.
\end{summarybox}

%%%%%%%%%%%%%%%%%%%%%%%%%%%%%%%%%%%%%%%%%%%%%%%%%%%%%%%%%%%%%%%%%%%%%%%%%%%%%%
\chapter{Online Applications}
%%%%%%%%%%%%%%%%%%%%%%%%%%%%%%%%%%%%%%%%%%%%%%%%%%%%%%%%%%%%%%%%%%%%%%%%%%%%%%

\syllabusstrip{Online applications.}

\begin{learningoutcomes}
\begin{itemize}
\item The eight reasons apps are effective marketing tools.
\item The four app business models.
\item The three rules of app creation strategy.
\item The eleven features an app should incorporate.
\item Designing a mobile website, and why responsive design\ix{responsive design} superseded it.
\item Mobile analytics --- its significance and the seven metric groups.
\end{itemize}
\end{learningoutcomes}

\section{Why Apps Matter for Marketing}

\begin{keybox}[High Yield]
``How do online applications serve as effective tools for mobile marketing? What
are the key features that businesses should incorporate into their apps to ensure
they are maximizing marketing potential?''
\pyq{Improvement CIE, 28 Aug 2024 --- 10 M}
\end{keybox}

\begin{itemize}
\item \sfb{Permanent presence} --- the app icon sits on the home screen, creating
      continuous brand visibility.
\item \sfb{Push notification access} --- a direct, free, owned communication
      channel that does not depend on an algorithm.
\item \sfb{Rich behavioural data} --- in-app analytics reveal exactly how each
      user behaves, enabling deep personalisation.
\item \sfb{Higher engagement and conversion} --- app users typically browse more,
      convert more and spend more than mobile-web users.
\item \sfb{Loyalty and habit formation} --- points, streaks and saved preferences
      increase switching cost.
\item \sfb{Device capability access} --- camera, GPS, biometrics, wallet and
      contacts enable experiences the web cannot match.
\item \sfb{Offline availability} --- content and coupons remain usable without a
      connection.
\item \sfb{Frictionless repeat purchase} --- saved payment details and one-tap
      reorder.
\end{itemize}

\section{App Business Models}

\begin{itemize}
\item Free apps supported by advertising.
\item Paid apps supported by download fees.
\item Premium apps supported by in-app commerce.
\item Free apps supported by brands interested in connecting with customers.
\end{itemize}

\section{App Creation Strategy --- Three Rules}

\begin{frameworkbox}[The Three Rules]
\textbf{1. Make sure your app solves a problem.} The most effective apps solve
some problem for the user --- they facilitate a purchase, provide content, create
brand preference, or some combination. Analyse which of these your app will solve
and begin the design process from there.

\textbf{2. Get inside the mind of your user.} Understand how the user will engage
with the app --- will they be at home or in the office? Do they want information,
or simply to engage with the brand?

\textbf{3. Design with the end in mind.} Is the purpose to create brand
preference? To facilitate a financial transaction? To reduce customer churn\ix{churn}? Or
all of these?
\end{frameworkbox}

\section{Key Features an App Should Incorporate}

\begin{itemize}
\item \sfb{Fast onboarding} --- value visible before any sign-up demand.
\item \sfb{Personalised home screen} driven by behaviour.
\item \sfb{Search and filtering} that works.
\item \sfb{Push notification preferences} the user controls.
\item \sfb{In-app offers, loyalty and wallet.}
\item \sfb{One-tap checkout} with saved payment methods.
\item \sfb{Order tracking and history.}
\item \sfb{In-app support and chat.}
\item \sfb{Social sharing and referral.}
\item \sfb{Offline mode} for core content.
\item \sfb{Analytics instrumentation} on every meaningful event.
\end{itemize}

\section{Creating a Mobile Website}

With a smartphone you can browse any website available on a PC, but those sites
are generally unsuitable for a phone because they were not designed for one. The
screen is much smaller and behaviour differs --- a user may browse and read for
hours on a desktop, but on a phone will typically read only a small amount of
information over short periods. How the information is displayed, and how much of
it there is, must therefore be reconsidered.

\noindent When designing a mobile website, consider:

\begin{itemize}
\item You can only view one screen at a time --- design the navigation
      accordingly.
\item There is not much room for text, so don't use much.
\item Use large buttons for key calls to action.
\item Think about font usage --- make sure the important content really stands
      out.
\end{itemize}

\begin{keybox}[Why Responsive Design Superseded the Separate Mobile Site]
The modern equivalent of this advice is \keytermas{responsive
design}{responsive design} --- one codebase that adapts fluidly to every screen.
It is now preferred over maintaining a separate mobile site, because it avoids
duplicate content, splits no ranking authority and halves maintenance.
\end{keybox}

\section{Mobile Analytics}

\begin{keybox}[High Yield]
``Explain the significance of mobile analytics in mobile marketing campaigns. What
key metrics should marketers track to measure the success of their mobile
strategies?''
\pyq{SEE, Oct/Nov 2023, Q8b --- 8 M}
\end{keybox}

\begin{exambox}[Significance]
Mobile analytics is the only way to see behaviour \emph{inside} an app, where web
analytics cannot reach. It reveals which acquisition sources deliver users who
actually retain rather than merely install; it exposes where users abandon the
onboarding or checkout flow; it enables cohort analysis so that retention can be
compared across releases; it supports A/B testing\ix{A/B testing} of in-app experiences; and it
links marketing spend to lifetime value rather than to installs.
\end{exambox}

\begin{figure}[htbp]
\centering
\begin{tikzpicture}[font=\sffamily\footnotesize,
    grp/.style={draw=ocre,line width=0.85pt,fill=ocrepale,rounded corners=3pt,
                minimum height=1.25cm,text width=1.72cm,align=center,
                font=\sffamily\bfseries\scriptsize,text=slate,inner sep=3pt}]
  \node[grp] (g1) at (0,0)     {ACQUISITION};
  \node[grp] (g2) at (2.0,0)   {ACTIVATION};
  \node[grp] (g3) at (4.0,0)   {ENGAGEMENT};
  \node[grp] (g4) at (6.0,0)   {RETENTION};
  \node[grp,draw=greenok,fill=greenpale] (g5) at (8.0,0) {MONETISATION};
  \foreach \a/\b in {g1/g2,g2/g3,g3/g4,g4/g5}{
    \draw[-{Latex[length=1.9mm]},ocre,line width=0.85pt] (\a)--(\b);}

  \node[grp,draw=deepblue,fill=deepbluepale] (g6) at (10.15,0.75)
    {NOTIFICATION};
  \node[grp,draw=deepblue,fill=deepbluepale] (g7) at (10.15,-0.75)
    {TECHNICAL};
  \node[anchor=north,text=slatelight,font=\sffamily\scriptsize\itshape,
        align=center,text width=11cm] at (5.1,-1.0)
    {the five funnel groups run left to right; notification and technical health
     cut across all of them};
\end{tikzpicture}
\caption[The mobile analytics metric groups]%
        {Five funnel groups plus two cross-cutting ones. Judging a campaign on
         installs alone stops at the first box.}
\label{fig:mobanalytics}
\end{figure}

\begin{mmlongtable}{The seven mobile analytics metric groups}{tab:mobmetrics}%
{@{}P{0.22\textwidth}P{0.71\textwidth}@{}}
\rowcolor{ocre}\thd{Metric group} & \thd{Metrics}\\
\midrule
\endfirsthead
\rowcolor{ocre}\thd{Metric group} & \thd{Metrics}\\
\midrule
\endhead
\rlab{Acquisition} & Installs, cost per install, install source and attribution,
click-to-install rate\\
\rowcolor{tablealt}
\rlab{Activation} & Onboarding completion rate, time to first key action,
registration rate\\
\rlab{Engagement} & Daily and monthly active users, stickiness (daily divided by
monthly actives), session length, sessions per user, screens per session\\
\rowcolor{tablealt}
\rlab{Retention} & Day-1, Day-7 and Day-30 retention, churn rate, cohort
retention curves\\
\rlab{Monetisation} & Conversion rate, average revenue per user, average order
value, lifetime value, the ratio of lifetime value to customer acquisition cost\\
\rowcolor{tablealt}
\rlab{Notification performance} & Push opt-in rate, delivery rate, open rate,
opt-out rate\\
\rlab{Technical} & Crash rate, app load time, API latency, app store rating\\
\end{mmlongtable}

\begin{summarybox}
Apps are effective marketing tools because they give permanent home-screen
presence, owned push access, rich behavioural data, higher engagement and
conversion, loyalty and habit formation, device capability access, offline
availability and frictionless repeat purchase. Four business models fund them.
Three rules govern their creation: solve a real problem, understand how the user
will engage, and design with the end purpose in mind. Eleven features maximise
marketing potential, from fast onboarding to full analytics instrumentation. A
mobile site must be designed for one screen at a time, little text and large
buttons --- and responsive design has superseded the separate mobile site because
it avoids duplicate content and split authority. Mobile analytics is the only view
inside the app, and it is measured in seven groups: acquisition, activation,
engagement, retention, monetisation, notification performance and technical
health.
\end{summarybox}
